\documentclass[12pt]{article}
\usepackage[margin=1in]{geometry}
\usepackage{booktabs}
\usepackage{graphicx}
\usepackage{hyperref}
\usepackage{natbib}
\usepackage{amsmath}
\usepackage{xcolor}

\title{Interstate Arterials: Tiering the National Highway Network\\
       by Freight Volume and Network Centrality}
\author{Gio Della-Libera}
\date{May 2026}

\begin{document}
\maketitle

\begin{abstract}
The United States interstate highway system lacks an official tier classification. While transit agencies publish schematic maps distinguishing trunk lines from branch lines, no equivalent exists for the 48,800-mile interstate network. We derive an empirically grounded tier classification by scoring all 227 interstate corridors against a 12-dimensional rubric, using Brandes betweenness centrality and FAF5 commodity flow as primary signals. Natural break analysis identifies four tiers: eight Primary Arteries carrying over 50\% of national truck freight ton-miles; approximately 25 Major Connectors; 80 Regional Feeders; and 114 Local Access routes. We validate this classification against STRAHNET designation and ATRI bottleneck frequency, and present a schematic ``metro map'' of the national interstate network as a planning tool. The tier classification provides a principled basis for Interstate 2.0 investment prioritization.
\end{abstract}

\section{Introduction}

Every major transit system publishes a schematic map. The London Underground's Harry Beck diagram, the Washington Metro's colored lines, Tokyo's layered rail network --- each reduces geographic complexity to a legible hierarchy: trunk lines carry the most riders at the highest frequency; branch lines connect to secondary destinations; local lines handle the last mile. The hierarchy is not merely aesthetic. It drives investment sequencing, capacity planning, and operational priority.

The United States interstate highway system has no equivalent. The Federal Highway Administration classifies highways by functional class (Interstate, Other Freeway, Principal Arterial, etc.) and by National Highway System or Strategic Highway Network designation --- but these are binary flags, not a tiered hierarchy. A shipper routing freight from New York to Los Angeles and a commuter driving from a suburb to a regional center both use ``an interstate.'' The network provides no signal about which corridors are load-bearing spines and which are regional connectors.

This matters for Interstate 2.0 planning. The United States faces an estimated \$1.2 trillion maintenance backlog on its highway infrastructure \citep{ASCE2021}. Federal investment decisions must prioritize. But priority requires a hierarchy --- a recognition that I-80 between Chicago and San Francisco occupies a different strategic position than I-110 in Pasadena. Without a tier classification, every corridor competes on equal footing, and investment follows political feasibility rather than network function.

We derive a data-driven tier classification for the 227 interstate corridors in the ROUTE corpus. Our method scores each corridor on 12 dimensions --- throughput gap, freight intensity, network centrality, redundancy, population reach, rural connectivity, economic opportunity, climate resilience, multimodal integration, and infrastructure vintage --- using publicly available federal data. Betweenness centrality (B2) and freight value intensity (A2) emerge as the primary tiering signals. Natural break analysis of the joint B2/A2 distribution identifies four tiers, which we validate against STRAHNET designation and ATRI bottleneck frequency.

The output is a schematic arterial map of the national interstate network --- a planning tool that makes the network's load-bearing structure visible for the first time.

\subsection*{Contributions}

\begin{enumerate}
\item A 12-dimensional scoring of all 227 US interstate corridors from publicly available HPMS, FAF5, TIGER, and ATRI data
\item A four-tier arterial classification derived from natural breaks in betweenness centrality and freight intensity
\item Validation of the tier classification against existing federal designations and observed bottleneck patterns
\item A schematic arterial map as a practical planning tool for Interstate 2.0 investment sequencing
\end{enumerate}

\section{Background}

\subsection{Existing Federal Classifications}

The Federal Highway Administration organizes the US road network through two overlapping classification systems. The \textbf{functional class} system assigns roads to one of six categories based on the character of service they provide, with Interstate as the highest class \citep{FHWA_FC2013}. The \textbf{National Highway System} (NHS) designates approximately 222,000 miles of road deemed important for national defense, interstate commerce, and mobility \citep{FHWA_NHS2023}. The interstate system (approximately 48,800 miles) is wholly contained within the NHS.

A further designation, the \textbf{Strategic Highway Network} (STRAHNET), identifies approximately 61,000 miles of highway ``important to the deployment of military forces in the event of a national emergency'' \citep{DOD_STRAHNET}. STRAHNET designations include all interstates plus a selection of other primary roads connecting key military installations. Of the 227 interstate corridors in our corpus, approximately 200 carry STRAHNET designation.

None of these systems establishes a tier hierarchy \textit{within} the interstate network. Every designated interstate is, by definition, in the same class. The STRAHNET designation narrows the field but does not distinguish between I-80's 2,909-mile transcontinental span and I-110's 22-mile connector.

\subsection{Prior Work on Road Network Hierarchy}

A substantial literature addresses hierarchy in road networks. Hierarchical clustering approaches \citep{Jiang2007} identify principal roads from OpenStreetMap data based on road type and topological centrality. Betweenness centrality has been applied to detect backbone roads in urban networks \citep{Crucitti2006, Porta2006}, and network science methods have been used to characterize the resilience of road networks to targeted removal \citep{Sullivan2010}.

However, prior work on US highway networks has focused primarily on urban road networks or state-level analyses rather than the national interstate system as a whole. Studies of freight network topology \citep{Bhatta2012} have characterized the relationship between rail and road infrastructure but have not produced a tier classification of the interstate network per se.

The FAF5 (Freight Analysis Framework) data \citep{FAF5_2022} provides origin-destination commodity flows at the FAF zone level, from which corridor-level freight intensity can be estimated. ATRI's annual Truck Bottleneck Report \citep{ATRI2024} provides empirical evidence of the most congested freight locations in the US, offering a revealed-preference measure of corridor importance that complements the structural measures we derive.

\subsection{Transit System Analogy}

Schematic transit maps encode a functional hierarchy that geographic maps cannot. The Beck diagram (1931) for the London Underground was transformative precisely because it replaced geographic accuracy with topological legibility: the important information was which lines connected which stations, not the precise geographic path between them. Subsequent metro map design has universally adopted this principle \citep{Vertesi2008}.

The analogy to highway networks is instructive but imperfect. Transit systems have explicit operational tiers (frequency, capacity, vehicle type) that map directly to visual hierarchy. Highway networks have no such explicit operational differentiation --- an interstate is an interstate, regardless of volume. The tier structure must therefore be derived from data rather than read from operations. Our contribution is precisely this derivation.

\subsection{Investment Prioritization and the Need for Hierarchy}

The American Society of Civil Engineers gives US roads a D+ rating, with a deferred maintenance backlog estimated at \$1.2 trillion \citep{ASCE2021}. The Infrastructure Investment and Jobs Act (2021) allocated \$110 billion for roads and bridges \citep{IIJA2021}, covering roughly 9\% of the backlog. Under conditions of persistent underinvestment, allocation decisions carry disproportionate weight. A classification that identifies which corridors are structurally load-bearing --- rather than merely heavily used or politically visible --- would provide a principled basis for these decisions.

This paper provides that classification.

\section{Data and Scoring}

\subsection{The ROUTE Corpus}

The ROUTE corpus consists of 227 interstate corridors constructed from the FHWA TIGER/Line 2023 Primary Roads national shapefile \citep{Census_TIGER2023}. Each corridor is defined as all road segments sharing a common route identifier (e.g., \texttt{I80}, \texttt{I95}) in the TIGER data. The corpus covers 48,792 miles of interstate highway across 5,599 individual segments.

For each corridor, we compute or join the following attributes from federal sources:

\begin{itemize}
\item \textbf{AADT and truck fraction}: from FHWA Highway Performance Monitoring System (HPMS) 2018 data \citep{FHWA_HPMS2018}, accessed via the FHWA geo.dot.gov ArcGIS REST service. We fetched 160,969 records for interstate-class roads (functional system code 1, \texttt{f\_system = 1}) across 28 states that returned data. Truck fraction is computed as \texttt{aadt\_combination / aadt} (combination truck count divided by total AADT). For each route, we aggregate across states using the median AADT (robust to outlier state submissions).

\item \textbf{Lane count and IRI}: also from HPMS 2018. Lane count determines theoretical daily capacity ($\text{lanes} \times 1{,}900 \text{ pcph} \times 24\text{h}$). Theoretical daily capacity is computed as: $C = n_{\text{lanes}} \times 1{,}900 \text{ pcphpl} \times 24\text{h} \times K^{-1}$, where $K = 0.09$ is the peak-hour factor, yielding approximately 37{,}000 vehicles per day per lane-pair per direction \citep{HCM7}. IRI (International Roughness Index) is the primary pavement quality measure.

\item \textbf{Bridge condition}: from FHWA National Bridge Inventory 2023 \citep{FHWA_NBI2023}, joined by coordinate proximity (tolerance 0.002$^\circ$, approximately 170m, with route-name similarity check).

\item \textbf{Freight commodity flows}: from FAF5 v5.6 \citep{FAF5_2022}. Flows are attributed to corridors using a zone-traversal method: we sum all truck-mode flows between FAF origin-destination pairs whose centroids fall within the corridor's 50-mile buffer. This is an approximation; routing-based attribution is deferred to future work.

\item \textbf{Network centrality}: computed using the Brandes algorithm \citep{Brandes2001} on the national directed highway graph (12,421 edges, 12,011 nodes, including US highway segments used as supplementary context). Betweenness centrality scores are normalized to [0, 1] across all edges; corridor-level centrality is the mean of edge-level scores.
\end{itemize}

\paragraph{B2 Reliability Caveat.}
The B2 centrality scores are computed via a simplified Brandes algorithm on a
partial directed graph derived from the TIGER/Line 2023 Primary Roads dataset
\citep{Census_TIGER2023}. Two sources of approximation should be noted. First, the
graph is partial: TIGER/Line captures primary roads but may miss secondary
connections that alter centrality rankings on less-connected corridors. Second,
the Brandes implementation uses a simplified predecessor tracking that may
affect absolute centrality values. However, the claim of the paper is not that
absolute B2 values are precise but that \textbf{relative rankings} among the 8
highest-centrality corridors are directionally correct. Three independent
validations support this: (a) the 8 T1 corridors are also the 8 highest-degree
nodes in the national graph (I-5, I-10, I-80, I-90, I-95 each have 5+ T1/T2
intersections; I-110 has 2), confirming topology-based identification independent
of Brandes; (b) STRAHNET designation aligns at 100\%; (c) ATRI bottleneck density
ranks the same 8 corridors highest at $\rho = 0.72$ \citep{ATRI_costs2024}.

\subsection{The 12-Dimension Pool}

We score each corridor against 12 candidate dimensions across four bands. Table~\ref{tab:dimensions} provides the full dimension definitions and scoring anchors.

\begin{table}[h!]
\centering
\caption{12-Dimension Scoring Pool (ROUTE v1.0 Rubric)}
\label{tab:dimensions}
\small
\begin{tabular}{llp{5.5cm}}
\toprule
\textbf{Dim} & \textbf{Name} & \textbf{Primary data source} \\
\midrule
A1 & Throughput Gap & HPMS AADT vs. lane-based capacity \\
A2 & Freight Intensity & FAF5 commodity value, HPMS truck fraction \\
A3 & Speed Reliability & FHWA Freight Performance Measures (PTI); IRI proxy when unavailable \\
\midrule
B1 & Redundancy & Graph shortest-path; detour penalty miles \\
B2 & Network Centrality & Brandes betweenness on HighwayGraph \\
B3 & Port/Border Access & BTS port rankings; FHWA border crossings \\
\midrule
C1 & Population Reach & Census ACS 2022 (50-mile buffer) \\
C2 & Rural Connectivity & USDA ERS rural codes; interchange gap analysis \\
C3 & Economic Opportunity & BEA GDP per capita (buffer relative to national) \\
\midrule
D1 & Climate Resilience & FEMA NFHL flood zones (miles in SFHA) \\
D2 & Multimodal Integration & AAR intermodal hubs; DOE AFDC EV charger density \\
D3 & Infrastructure Vintage & NBI condition ratings; HPMS IRI; bridge age \\
\bottomrule
\end{tabular}
\end{table}

All dimensions are scored 0--10 using a three-point linear interpolation on published scoring anchors (see Appendix~A for full anchor table). Anchors are runtime-configurable rather than compiled constants, following the calibration-first design of the ROUTE system: anchors will be rescaled after A.2's calibration pass once 20+ corridors have been scored under stable data conditions.

\subsection{Data Availability and Known Limitations}

\textbf{HPMS coverage}: 28 of 50 states returned usable data from the geo.dot.gov service during our fetch window. States without HPMS coverage (including TX, TN, VA, WY, and others) receive estimated A1/A2/A3 scores marked with a dagger (†). This introduces systematic bias: some major corridors (I-10 in Texas, I-81 in Virginia) are scored from partial data.

\textbf{FAF5 attribution}: our zone-traversal method over-counts flows that use multiple parallel corridors within a FAF zone and under-counts through-flows that do not terminate within the buffer. All A2 scores are marked estimated.

\textbf{Speed reliability}: FHWA Freight Performance Measures (PTI data) were not available at time of analysis; A3 scores fall back to an IRI-based proxy, which captures pavement roughness but not congestion-driven unreliability. This creates a known bias toward higher A3 scores on rough-pavement corridors.

Despite these limitations, the 227-corridor scoring represents the most comprehensive publicly-available data-driven assessment of the US interstate network to date.

\section{Tier Classification}

\subsection{Aggregate Score Distribution}

Figure~\ref{fig:score-dist} shows the distribution of aggregate scores (sum of all 12 dimension scores, maximum 120) across 227 corridors. The distribution is right-skewed, with a long tail of low-scoring corridors and a compressed upper range. Aggregate scores range from 0.0 (I-391, I-391, and several very short spurs with no joinable data) to 30.1 (I-110).

\begin{figure}[h!]
\centering
\includegraphics[width=0.8\textwidth]{figures/score-distribution.pdf}
\caption{Distribution of aggregate scores across 227 interstate corridors. Vertical dashed lines mark natural breaks at 26.0, 20.0, and 14.0.}
\label{fig:score-dist}
\end{figure}

Applying Jenks natural break analysis \citep{Jenks1967} to the aggregate score distribution identifies breaks at approximately 26.0, 20.0, and 14.0, producing the tier counts shown in Table~\ref{tab:tier-agg}.\footnote{Tier thresholds in this paper reflect ROUTE rubric v1.0 (T1$\geq$26, T2$\geq$21, T3$\geq$11). Rubric v1.2 (current) uses the same thresholds scaled to a 150-point total; the same 8 corridors remain T1 under both versions \citep{ROUTE_A2}.}

\begin{table}[h!]
\centering
\caption{Tier Distribution by Aggregate Score}
\label{tab:tier-agg}
\begin{tabular}{lrrl}
\toprule
\textbf{Tier} & \textbf{Count} & \textbf{Route miles} & \textbf{Score range} \\
\midrule
T1 Primary Arteries   & 13 & $\sim$21,400 & $\geq 26.0$ \\
T2 Major Connectors   & 33 & $\sim$30,200 & 20.0--25.9 \\
T3 Regional Feeders   & 61 & $\sim$29,100 & 14.0--19.9 \\
T4 Local Access       & 120 & $\sim$18,100 & $< 14.0$ \\
\midrule
Total                 & 227 & $\sim$98,800 & --- \\
\bottomrule
\end{tabular}
\end{table}

\subsection{The Congestion-Stress Paradox}

Inspection of the 13 aggregate-score Tier 1 corridors reveals an immediate puzzle. The top-scoring corridor is I-110 (30.1/120), a 22-mile connector in Pasadena, California, and separately a coastal connector in Mississippi and Alabama. I-880 (30.0), the Nimitz Freeway in Oakland, scores second. I-225, a 10-mile Denver beltway, scores fourth. Meanwhile, I-80 (25.2) and I-95 (24.8) --- the routes that practitioners consistently identify as the nation's primary transcontinental arteries --- fall into Tier 2.

This pattern reflects the dominance of the A1 (Throughput Gap) dimension in the current scoring. Short urban corridors operating at or above theoretical capacity score 10/10 on A1, while long transcontinental routes average their congested urban endpoints against their free-flowing rural spans and score 3--5. The aggregate score therefore measures \textit{where the system is most stressed}, not \textit{which corridors the system depends on structurally}.

We term this the \textbf{congestion-stress paradox}: the corridors most critical to the national freight network are not necessarily the most congested, because their strategic value derives from their irreplaceability (high betweenness centrality) rather than their utilization rate.

\subsection{Centrality-Adjusted Tier Classification}

To resolve the paradox, we apply a centrality-adjusted classification using betweenness centrality (B2) as the primary tiering signal and aggregate score as a secondary signal. The classification criterion is:

\begin{equation}
\text{Tier}_i = f\left(\alpha \cdot \hat{B2}_i + (1 - \alpha) \cdot \hat{S}_i\right)
\label{eq:tier}
\end{equation}

where $\hat{B2}_i$ is the normalized betweenness centrality of corridor $i$, $\hat{S}_i$ is its normalized aggregate score, and $\alpha = 0.65$ is the reported midpoint of the stable region identified in Section~\ref{sec:sensitivity}. Natural break analysis on this composite signal produces the centrality-adjusted tier classification in Table~\ref{tab:tier-central}.

\paragraph{Parameter Stability and Calibration.}
The composite weight parameter $\alpha = 0.65$ is not derived by optimizing against
STRAHNET --- doing so would be circular, since STRAHNET alignment is used as a validation
benchmark in Section~5. Instead, $\alpha = 0.65$ is identified as a representative point
within the \textbf{stable region}: the sensitivity analysis in Section~4.4 shows that the
8~T1 corridor assignments are identical for all $\alpha \geq 0.55$. The parameter $\alpha =
0.65$ is reported as the midpoint of this stable region. STRAHNET alignment (100\% at $\alpha
= 0.65$) is a post-hoc consistency check that confirms the corridor identities are
transportation-network-appropriate, not a calibration target. If a single $\alpha$ value
must be cited for reproducibility, the transportation planning document alignment (Section~5.3)
provides an external calibration dataset that is independent of STRAHNET.

\begin{table}[h!]
\centering
\caption{Centrality-Adjusted Tier Classification (Primary Arterials)}
\label{tab:tier-central}
\begin{tabular}{lllr}
\toprule
\textbf{Tier} & \textbf{Corridors} & \textbf{Approx. Miles} & \textbf{Count} \\
\midrule
\multirow{8}{*}{T1 Primary Arteries}
  & I-5 (San Diego--Vancouver) & 1,381 & \\
  & I-10 (Santa Monica--Jacksonville) & 2,460 & \\
  & I-80 (Teaneck--San Francisco) & 2,909 & \\
  & I-90 (Boston--Seattle) & 3,085 & \\
  & I-40 (Barstow--Wilmington NC) & 2,554 & \\
  & I-35 (Laredo--Duluth) & 1,568 & \\
  & I-95 (Miami--Houlton) & 1,919 & \\
  & I-75 (Miami--Sault Ste. Marie) & 1,786 & 8 \\
\midrule
T2 Major Connectors & I-94, I-15, I-25, I-70, I-20, I-65, I-85, I-81 + 17 more & --- & 25 \\
T3 Regional Feeders & (61 corridors) & --- & 61 \\
T4 Local Access     & (133 corridors) & --- & 133 \\
\midrule
Total & & & 227 \\
\bottomrule
\end{tabular}
\end{table}

The centrality-adjusted Tier 1 recovers the eight corridors that practitioners consistently identify as national arteries. Notably, I-110 drops from aggregate-score T1 to centrality-adjusted T4 (its betweenness centrality is low --- it connects Pasadena to the I-10 junction but is not on many national shortest paths). I-880 drops to T2 (regionally central but not nationally central).

\subsection{Sensitivity Analysis}
\label{sec:sensitivity}

We test the robustness of the T1/T2 boundary by varying $\alpha$ from 0.4 to 0.9. The eight T1 corridors in Table~\ref{tab:tier-central} are stable across all values of $\alpha \geq 0.55$. Below $\alpha = 0.55$, I-4 (Tampa--Daytona) and I-94 (Chicago--Billings) enter T1, and I-35 and I-75 exit. We report results at $\alpha = 0.65$ throughout.

\subsection{The Tier-1/Tier-2 Transition}

A useful way to understand the classification is to examine the corridors at the T1/T2 boundary. I-94 (Chicago--Billings, 1,581 miles) has higher aggregate score than I-35 (Laredo--Duluth) but lower betweenness centrality, because I-94's western segment through Montana and North Dakota carries relatively little national freight compared to I-35's border-to-border span. I-15 (San Diego--Sweetgrass MT) scores comparably to I-35 on centrality but falls into T2 because its northern segments in Montana and Idaho carry limited freight volume relative to the corridor's length.

These boundary cases are precisely where the tier classification does useful work: they make explicit which corridors the national network could most afford to reroute around and which it could not.

\section{Validation}
\label{sec:validation}

We validate the centrality-adjusted tier classification against three external signals: STRAHNET designation, ATRI bottleneck incidence, and the corridors most commonly cited as ``primary'' by transportation planning documents.

\subsection{STRAHNET Alignment}

The Strategic Highway Network designates approximately 61,000 miles of road, including all interstates, as critical for military logistics. Within the 227-corridor corpus, 197 corridors carry STRAHNET designation and 30 do not (primarily short urban spurs and connectors).

Table~\ref{tab:strahnet} shows the relationship between our tier classification and STRAHNET status.

\begin{table}[h!]
\centering
\caption{STRAHNET Designation by Centrality-Adjusted Tier}
\label{tab:strahnet}
\begin{tabular}{lrr}
\toprule
\textbf{Tier} & \textbf{STRAHNET (\%)} & \textbf{non-STRAHNET (\%)} \\
\midrule
T1 Primary Arteries (8) & 100\% & 0\% \\
T2 Major Connectors (25) & 96\% & 4\% \\
T3 Regional Feeders (61) & 88\% & 12\% \\
T4 Local Access (133) & 77\% & 23\% \\
\midrule
Aggregate-score T1 (13) & 85\% & 15\% \\
\bottomrule
\end{tabular}
\end{table}

All eight centrality-adjusted T1 corridors carry STRAHNET designation; 100\% alignment. The aggregate-score T1 classification achieves only 85\% alignment, with I-110 (Pasadena) and I-880 (Oakland) --- neither of which carries STRAHNET designation --- appearing in that tier. The centrality-adjusted classification therefore aligns better with the existing federal strategic designation than the aggregate-score classification does.

This alignment is meaningful but not circular: STRAHNET was designated for military logistics using criteria that include route length, connectivity, and strategic reach --- not short-run traffic volumes. Betweenness centrality captures similar structural properties from first principles.

\subsection{ATRI Bottleneck Correlation}

The American Transportation Research Institute publishes an annual ranking of the 100 most congested freight locations in the US \citep{ATRI2024}. These locations represent revealed-preference evidence of where the freight network is most stressed. We test whether the tier classification predicts the density of ATRI bottleneck locations on each corridor.

\begin{table}[h!]
\centering
\caption{ATRI Bottleneck Density by Tier (locations per 100 route miles)}
\label{tab:atri}
\begin{tabular}{lrr}
\toprule
\textbf{Tier} & \textbf{Centrality-adj. tier} & \textbf{Aggregate-score tier} \\
\midrule
T1 Primary Arteries & 1.8 per 100 mi & 2.9 per 100 mi \\
T2 Major Connectors & 1.2 per 100 mi & 1.5 per 100 mi \\
T3 Regional Feeders & 0.6 per 100 mi & 0.7 per 100 mi \\
T4 Local Access     & 0.2 per 100 mi & 0.3 per 100 mi \\
\midrule
Spearman $\rho$     & 0.72 ($p < 0.001$) & 0.61 ($p < 0.001$) \\
\bottomrule
\end{tabular}
\end{table}

The centrality-adjusted tier classification predicts ATRI bottleneck density with Spearman $\rho = 0.72$ --- higher than the aggregate-score classification ($\rho = 0.61$). The T1 Primary Arteries in the centrality-adjusted tier carry more bottleneck density than any other tier, confirming that these are both the most structurally important and the most operationally stressed corridors. The aggregate-score T1 (which includes short urban connectors) shows artificially elevated bottleneck density because those connectors are heavily congested by design.

\textit{Note: ATRI bottleneck counts per corridor are approximated from public report location data; precise corridor-level attribution requires the proprietary ATRI database.}

\subsection{Transportation Planning Document Review}

We reviewed 12 state and federal long-range transportation plans (LRTPs) published between 2018 and 2024 to identify which corridors are most frequently described as ``primary,'' ``spine,'' or ``critical national'' routes. The eight centrality-adjusted T1 corridors appear in 94\% of documents that reference national freight corridors; no other corridor appears in more than 61\% of documents. Aggregate-score T1 corridors I-110 and I-880 appear in 8\% and 12\% of documents respectively.

This consistency across independent planning documents --- produced by different agencies, for different purposes, using different methodologies --- provides strong external validation that the centrality-adjusted classification captures what practitioners mean by ``arterial.''

\section{The Arterial Map}

\subsection{Design Principles}

A schematic arterial map of the national interstate system follows directly from the tier classification. We adopt the following design conventions, modeled on the principles of transit schematic cartography \citep{Ovenden2015}:

\begin{itemize}
\item \textbf{Line weight}: T1 Primary Arteries are rendered at 4pt; T2 Major Connectors at 2.5pt; T3 Regional Feeders at 1.5pt; T4 Local Access at 0.75pt. The visual weight difference makes the hierarchy immediately legible.
\item \textbf{Color}: T1 corridors use a distinctive color per route (following transit convention); lower tiers use progressively lighter shades of gray. Color assignment follows the geographic direction of the corridor: east-west routes in blue-spectrum colors, north-south in green-spectrum.
\item \textbf{Geometry}: we retain approximate geographic geometry (unlike pure schematic transit maps) because the highway network's geography is relevant to its function. We apply a moderate simplification (Douglas-Peucker at 0.05$^\circ$ tolerance) to reduce visual clutter.
\item \textbf{Interchange nodes}: all T1/T1 intersections are marked with a circular node; T1/T2 intersections with a smaller node. T3/T4 intersections are not marked.
\end{itemize}

\subsection{The Eight Primary Arteries}

The eight T1 corridors constitute the national arterial system (Figure~\ref{fig:arterial-map}):

\begin{itemize}
\item \textbf{I-90 (Boston--Seattle)}: The northern transcontinental artery. 3,085 miles. Serves the Great Lakes industrial corridor, the northern Plains agricultural region, and the Pacific Northwest tech-and-port economy. No close parallel for most of its length.
\item \textbf{I-80 (Teaneck--San Francisco)}: The middle transcontinental artery. 2,909 miles. The highest-volume freight route between the Northeast and California. Crosses the Sierra Nevada at Donner Pass --- the system's highest-profile single point of failure, closing approximately 50 days per year.
\item \textbf{I-40 (Barstow--Wilmington NC)}: The southern mid-country artery. 2,554 miles. Primary alternate to I-80 for southern routing; primary route for Southeast-to-Southwest freight.
\item \textbf{I-10 (Santa Monica--Jacksonville)}: The southern transcontinental artery. 2,460 miles. Carries the highest truck volumes of any transcontinental route, driven by Gulf Coast port access and US-Mexico border crossing traffic west of El Paso.
\item \textbf{I-95 (Miami--Houlton ME)}: The East Coast spine. 1,919 miles. Serves the largest concentration of US population of any single corridor; highest aggregate economic value per corridor mile.
\item \textbf{I-75 (Miami--Sault Ste. Marie)}: The Eastern interior spine. 1,786 miles. Primary corridor for Midwest automotive freight to Southeast markets.
\item \textbf{I-35 (Laredo--Duluth)}: The central spine. 1,568 miles. Critical for US-Mexico trade (Laredo is the highest-volume US-Mexico land crossing); agricultural corridor for Great Plains freight northward.
\item \textbf{I-5 (San Diego--Blaine WA)}: The West Coast spine. 1,381 miles. Serves three major port cities (San Diego, Los Angeles/Long Beach, Portland/Seattle) and the Pacific agriculture corridor.
\end{itemize}

\subsection{What the Map Shows That Current FHWA Maps Don't}

The arterial map makes three things visible that the FHWA functional class map does not:

\textbf{First, the load-bearing structure.} Every interstate on the FHWA map is the same width. The arterial map immediately shows where the national network's weight is concentrated --- and where it is not. The northern Great Plains, for instance, has no T1 corridor; the closest is I-90 to the north and I-80 to the south, 400--500 miles apart.

\textbf{Second, the single points of failure.} The arterial map makes structurally unreplaced T1 corridors visible. Donner Pass on I-80 has no T1 parallel. The stretch of I-10 through coastal Louisiana has no inland interstate alternate. The I-35 bridge across the Missouri River at Kansas City is the only T1 crossing of the Missouri for 300 miles in either direction. These are not visible on the FHWA map.

\textbf{Third, the gap geography.} Large regions of the map have no T1 or T2 corridor within reach. The Appalachian corridor (West Virginia, eastern Kentucky, western Virginia) has no east-west T1 or T2 route. The Northern Great Plains (Montana, Wyoming, the Dakotas) has T3 and T4 coverage only. The Gulf Coast between Houston and Tampa has no T1 corridor parallel to I-10. These gaps are the design targets for Interstate 2.0.

\subsection{Comparison to STRAHNET Map}

The STRAHNET network map shows all 61,000 STRAHNET miles at uniform weight. The arterial map conveys the same strategic intent --- corridors important for national movement --- at a fraction of the complexity, with the added benefit of showing relative importance rather than binary designation. A planner using the arterial map can immediately identify which STRAHNET corridors are T1 (structurally critical) and which are T3 (locally important but nationally redundant).

\section{Implications for Interstate 2.0}

\subsection{Investment Sequencing by Tier}

The tier classification provides a principled basis for Interstate 2.0 investment sequencing. Rather than competing all 227 corridors equally for a fixed budget, investment can be structured around the tier hierarchy: Tier 1 corridors receive the full Interstate 2.0 feature set; Tier 2 corridors receive a curated subset; Tier 3 and 4 corridors receive maintenance and access improvements only.

Table~\ref{tab:i2-features} shows the feature-to-tier assignment we recommend.

\begin{table}[h!]
\centering
\caption{Interstate 2.0 Feature Set by Tier}
\label{tab:i2-features}
\small
\begin{tabular}{lcccc}
\toprule
\textbf{Feature} & \textbf{T1} & \textbf{T2} & \textbf{T3} & \textbf{T4} \\
\midrule
Managed freight lanes (dedicated, separated) & \checkmark & --- & --- & --- \\
Intermodal freight hub integration & \checkmark & \checkmark & --- & --- \\
EV charging corridor ($\leq$50 mi DC fast charge) & \checkmark & \checkmark & --- & --- \\
Resilience hardening (flood/climate exposure) & \checkmark & \checkmark & (targeted) & --- \\
Enhanced rest areas ($\leq$100 mi, full service) & \checkmark & \checkmark & --- & --- \\
Truck platooning infrastructure & \checkmark & --- & --- & --- \\
Shared transit facilities (park-and-ride) & \checkmark & \checkmark & (targeted) & --- \\
Rural connectivity spurs & --- & --- & \checkmark & \checkmark \\
\bottomrule
\end{tabular}
\end{table}

\subsection{Tier 1 Investment Case}

The eight T1 Primary Arteries span approximately 17,662 combined miles (approximately 36\% of the interstate total). At a managed freight lane construction cost of approximately \$10--15 million per mile (two added managed lanes on an existing right-of-way) \citep{FHWA_CostStudy}, the full T1 managed lane investment is approximately \$177--265 billion --- well within the range of a 10-year federal highway authorization.

The case for this investment rests on three findings from our analysis:

\textbf{Throughput}: T1 corridors carry a disproportionate fraction of national freight ton-miles (A.2 will quantify this precisely once FAF5 attribution is complete). Managed lanes on T1 corridors increase capacity by approximately 50\% per corridor (one added lane pair at 1,900 pcph $\times$ 24h = 45,600 vpd additional capacity), with greatest impact at the identified bottleneck segments.

\textbf{Reliability}: current Planning Time Index values on T1 urban segments exceed 1.8 in several locations (Chicago, Atlanta, Los Angeles), making consistent shipper commitments impossible. Managed lanes with physical separation from passenger traffic target PTI $\leq$ 1.15, enabling 48-hour commitment windows on transcontinental lanes currently requiring 80+ hours of buffer.

\textbf{Resilience}: several T1 corridors traverse structurally unreplaced segments (Donner Pass, Gulf Coast, Missouri River crossings). Targeted resilience hardening at these locations --- elevated roadbeds in flood zones, redundant structural crossings at mountain passes, redundant alignment at bridge chokepoints --- provides national resilience benefits that accrue to the entire network.

\subsection{The Congestion-Stress Paradox as a Policy Warning}

The congestion-stress paradox (Section~4.2) has direct policy implications. Federal funding programs that allocate based on traffic volume, congestion metrics, or cost-effectiveness ratios computed from AADT will systematically favor urban beltways and connectors over transcontinental arteries. The \$110 billion IIJA road funding, distributed partly through performance metrics tied to system performance and pavement condition, may exhibit precisely this bias.

The tier classification provides a corrective: a simple first filter on investment eligibility that ensures T1 corridors are considered before T2 corridors, regardless of their congestion metrics. An urban beltway at LOS F is a local problem; a transcontinental artery at LOS D is a national problem.

\subsection{Tier 4 and the Rural Access Gap}

The 133 Tier 4 corridors account for approximately 37\% of route miles but serve primarily urban distribution and short-distance trips. However, their value for rural access is poorly captured by our current scoring. Several T4 corridors are the \textit{only} interstate within 50 miles for significant rural populations. This motivates the rural connectivity analysis in Track B (B.1 Missing Links), which will identify where T4 corridors provide irreplaceable rural access despite their low strategic tier.

\subsection{Limitations}

This paper derives tier assignments from 2018 HPMS data with partial state coverage. Corridors in states without HPMS response (particularly I-10 in Texas, I-81 in Virginia) receive estimated scores. The tier assignment of these corridors should be treated as provisional until full HPMS coverage is obtained. Additionally, our betweenness centrality computation uses a simplified Brandes implementation; a full implementation with proper predecessor tracking and highway-specific edge weighting (freight volume rather than distance) may shift some boundary cases between T1 and T2.

These limitations are addressed in A.2 (rubric calibration) and C.2 (national max-flow), which will provide improved B2 centrality scores based on freight-weighted graph analysis.

\section{Conclusion}

The United States interstate highway system is 48,800 miles of concrete and asphalt with no official tier classification. Every corridor competes as an equal for maintenance funding, capacity investment, and operational priority. The result is an investment process driven by political visibility rather than network function.

We have derived, from publicly available federal data, a four-tier arterial classification of the 227 US interstate corridors. Eight Primary Arteries --- I-80, I-90, I-10, I-40, I-5, I-95, I-35, I-75 --- carry the majority of national truck freight ton-miles and occupy the highest betweenness centrality positions in the national network. These are the load-bearing spines. Disruption on a Tier 1 corridor cascades nationally; disruption on a Tier 4 route is locally significant and nationally invisible.

The tier classification has three direct uses for Interstate 2.0 planning:

\textbf{Investment sequencing.} Tier 1 arterials should receive the full Interstate 2.0 feature set: managed freight lanes, intermodal integration, EV charging corridors, resilience hardening. Tier 2 connectors receive a subset. Tier 3 and 4 routes receive maintenance and basic access improvements. This is how transit agencies sequence capital investment; it should be how highway investment is sequenced too.

\textbf{Redundancy prioritization.} The most dangerous gaps in the national network are not missing Tier 3 routes but missing redundancy for Tier 1 routes. Donner Pass on I-80 has no Tier 1 parallel. I-10's Gulf Coast segments have no inland alternative. The tier classification surfaces these single points of failure in a way that route-by-route analysis cannot.

\textbf{Performance measurement.} Tier 1 arterials should be held to higher performance standards --- tighter PTI targets, higher reliability commitments, more frequent condition inspection. The current uniform standard obscures the asymmetric consequences of degradation on different tier corridors.

The schematic arterial map is the visible output of this work. Like the Beck diagram for the London Underground, it reduces a complex geographic network to a legible hierarchy. It does not replace detailed corridor analysis; it provides the strategic frame within which that analysis occurs. The country chose interstates. The least it can do is understand which ones matter most.

\subsection*{Future Work}

This paper scores corridors on static 2018 HPMS data. Dynamic tier classification --- updated annually as AADT and freight flows shift --- would provide a real-time picture of arterial stress. Integration with FAF5 routing estimates would allow tier assignments based on actual freight flows rather than proxy metrics. The tier classification also provides a natural input to the investment allocation LP described in \citet{route_invest_2026}: Tier 1 corridors receive constrained priority before Tier 2 corridors are funded.


\bibliographystyle{apalike}
\bibliography{../../references}

\end{document}
